\documentclass[11pt,a4paper]{article}

%================================================
% ENCODAGE ET LANGUE
%================================================
\usepackage[utf8]{inputenc}
\usepackage[T1]{fontenc}
\usepackage[french]{babel}

%================================================
% MISE EN PAGE
%================================================
\usepackage{geometry}
\geometry{margin=1.55cm}

%================================================
% PACKAGE GTOLI
%================================================
\usepackage{../styles/gtoli}

%================================================
% INFORMATIONS DU DOCUMENT
%================================================
\title{\textbf{Physique}\\[0.2cm]
\large Titre du chapitre}

\author{GToli}
\date{\today}

%================================================
% DOCUMENT
%================================================
\begin{document}

\maketitle

%================================================
\section{Notion physique concernée}
%================================================

\begin{cadreobjectif}
Dans ce chapitre, l’élève doit être capable de :
\begin{itemize}
    \item identifier les grandeurs physiques en jeu ;
    \item utiliser les unités correctes ;
    \item appliquer une loi physique ;
    \item interpréter un résultat ;
    \item représenter une situation par un schéma ;
    \item vérifier la cohérence physique d’une réponse.
\end{itemize}
\end{cadreobjectif}

%================================================
\section{Situation physique de départ}
%================================================

\begin{cadrebleu}{Mise en contexte}
Cette section permet d’introduire la situation physique étudiée.

\medskip

On précisera :
\begin{itemize}
    \item le phénomène observé ;
    \item les grandeurs mesurables ;
    \item les hypothèses éventuelles ;
    \item la question physique posée.
\end{itemize}
\end{cadrebleu}

%================================================
\section{Grandeurs, symboles et unités}
%================================================

\begin{cadregris}{Tableau des grandeurs}
\renewcommand{\arraystretch}{1.35}
\begin{tabularx}{\textwidth}{>{\bfseries}p{3cm}p{2.5cm}p{2.5cm}X}
\toprule
Grandeur & Symbole & Unité SI & Signification \\
\midrule
Grandeur 1 & $x$ & m & Déplacement ou distance \\
Grandeur 2 & $t$ & s & Durée \\
Grandeur 3 & $v$ & m/s & Vitesse \\
Grandeur 4 & $F$ & N & Force \\
\bottomrule
\end{tabularx}
\end{cadregris}

%================================================
\section{Rappel théorique}
%================================================

\begin{cadretheorie}
Cette section contient :
\begin{itemize}
    \item les définitions ;
    \item les lois physiques ;
    \item les formules ;
    \item les hypothèses d’application ;
    \item les conventions de signe ;
    \item les unités à utiliser.
\end{itemize}

\medskip

\textbf{Point essentiel :} en physique, une formule n’a de sens que si les grandeurs, les unités et le contexte sont correctement identifiés.
\end{cadretheorie}

%================================================
\section{Loi physique principale}
%================================================

\begin{cadrevert}{Loi ou relation fondamentale}
Écrire ici la relation principale du chapitre :

\[
\text{Formule à compléter}
\]

\medskip

\textbf{Signification des symboles :}
\begin{itemize}
    \item $x$ : grandeur à préciser ;
    \item $t$ : grandeur à préciser ;
    \item $v$ : grandeur à préciser.
\end{itemize}

\medskip

\textbf{Conditions d’application :}
\begin{itemize}
    \item condition 1 ;
    \item condition 2 ;
    \item condition 3.
\end{itemize}
\end{cadrevert}

%================================================
\section{Schéma de la situation}
%================================================

\begin{cadregris}{Représentation}
Insérer ici un schéma, une figure, un circuit, un diagramme ou une représentation simplifiée.

\medskip

\begin{center}
\begin{tikzpicture}
\draw[thick] (0,0) -- (6,0);
\draw[fill=gray!20] (1,0) rectangle (2,0.7);
\draw[->,thick] (2.1,0.35) -- (4,0.35) node[right] {$\vec{F}$};
\node at (1.5,1) {objet};
\end{tikzpicture}
\end{center}
\end{cadregris}

%================================================
\section{Protocole de résolution}
%================================================

\begin{cadremethode}
Pour résoudre un exercice de physique :

\begin{enumerate}
    \item lire attentivement la situation ;
    \item identifier le phénomène étudié ;
    \item relever les données ;
    \item convertir les unités si nécessaire ;
    \item choisir la loi physique adaptée ;
    \item isoler l’inconnue ;
    \item effectuer le calcul ;
    \item vérifier l’unité du résultat ;
    \item interpréter physiquement la réponse.
\end{enumerate}
\end{cadremethode}

%================================================
\section{Analyse dimensionnelle}
%================================================

\begin{cadreorange}{Vérification des unités}
Avant de conclure, il faut vérifier que l’unité obtenue correspond bien à la grandeur recherchée.

\medskip

Exemple :

\[
[v] = \frac{[x]}{[t]} = \frac{\mathrm{m}}{\mathrm{s}} = \mathrm{m/s}
\]

\medskip

Cette vérification permet de repérer rapidement certaines erreurs de formule ou de conversion.
\end{cadreorange}

%================================================
\section{Exemple modèle}
%================================================

\begin{cadreexemple}
\textbf{Énoncé :}

Insérer ici un exercice représentatif du chapitre.

\medskip

\textbf{Analyse physique :}
\begin{itemize}
    \item Quel phénomène est étudié ?
    \item Quelles grandeurs sont connues ?
    \item Quelle grandeur cherche-t-on ?
    \item Quelle loi physique faut-il appliquer ?
\end{itemize}

\medskip

\textbf{Résolution :}

Présenter les étapes avec unités et interprétation.
\end{cadreexemple}

%================================================
\section{Erreurs fréquentes}
%================================================

\begin{cadreattention}
Les erreurs fréquentes en physique sont :
\begin{itemize}
    \item oublier de convertir les unités ;
    \item utiliser une formule sans vérifier son domaine d’application ;
    \item confondre grandeur et unité ;
    \item oublier le sens physique du résultat ;
    \item négliger le schéma ;
    \item donner une réponse numérique sans interprétation.
\end{itemize}
\end{cadreattention}

%================================================
\section{Exercice d’application}
%================================================

\begin{cadreenonce}
\textbf{Exercice 1 :}

Insérer ici un exercice d’application directe.
\end{cadreenonce}

\begin{cadreaide}
Pour commencer :
\begin{itemize}
    \item repérer les données ;
    \item noter les unités ;
    \item convertir si nécessaire ;
    \item choisir la bonne loi ;
    \item isoler l’inconnue.
\end{itemize}
\end{cadreaide}

%================================================
\section{Solution détaillée}
%================================================

\begin{cadresolution}
La solution doit contenir :
\begin{enumerate}
    \item les données ;
    \item les conversions ;
    \item la loi utilisée ;
    \item l’isolation de l’inconnue ;
    \item le calcul ;
    \item l’unité ;
    \item l’interprétation physique.
\end{enumerate}
\end{cadresolution}

%================================================
\section{Activité expérimentale}
%================================================

\begin{cadreactivite}
Cette section peut contenir :
\begin{itemize}
    \item une manipulation ;
    \item un protocole expérimental ;
    \item un tableau de mesures ;
    \item une exploitation graphique ;
    \item une conclusion expérimentale.
\end{itemize}
\end{cadreactivite}

%================================================
\section{Synthèse}
%================================================

\begin{cadresynthese}
À retenir :

\begin{itemize}
    \item phénomène étudié : \dotfill
    \item grandeur principale : \dotfill
    \item loi physique : \dotfill
    \item unité importante : \dotfill
    \item erreur à éviter : \dotfill
    \item interprétation physique : \dotfill
\end{itemize}
\end{cadresynthese}

%================================================
\section{Auto-évaluation}
%================================================

\begin{cadregris}{Je vérifie ma maîtrise}
\begin{itemize}
    \item[$\square$] Je sais identifier les grandeurs.
    \item[$\square$] Je connais les unités.
    \item[$\square$] Je sais convertir.
    \item[$\square$] Je sais choisir la bonne loi.
    \item[$\square$] Je sais faire un schéma.
    \item[$\square$] Je sais interpréter le résultat.
\end{itemize}
\end{cadregris}

\end{document}